\chapter{Общепроектная проверка литературы и новизны}

Перед выбором архитектуры проверяются весь архив проекта и первичная
литература, чтобы не открыть повторно уже закрытый механизм и не объявить
известный приём собственной новизной.

\section{Область проверки}

Просмотрены архивные тексты RPFT/UGSM/TOE, Тома II--IV, реестр теорем,
итоговые заморозки, ретроспективные и обратные проверки, ветви носителя,
Вильсона--дефекта, поколений и Пати--Салама, а также соответствующие
машинные результаты. Внешняя сверка охватывает вариационные принципы Гиббса,
относительную энтропию, информационную геометрию, спектральное действие,
однопетлевые статистические суммы, Пати--Салам, механизм Хосотани,
вихри порядка три, фиксированный заряд, моды Майораны, конусы отображения,
детерминантные линии и четырёхфермионные каналы от кручения.

\section{Уровни новизны}

\begin{definition}[Уровни новизны]
\begin{enumerate}
 \item \(N_0\): известный элемент литературы;
 \item \(N_1\): его специализация к данным проекта;
 \item \(N_2\): ранее не полученное совместное утверждение уровня теоремы;
 \item \(N_3\): физическое замыкание со слепыми предсказаниями и независимым
 воспроизведением.
\end{enumerate}
\end{definition}

Само спектральное действие, мера Гиббса, алгебра Пати--Салама,
\(A_4\)-вихрь, вильсоновское нарушение симметрии, индекс дефекта Майораны
или гессиан логарифма определителя не являются новизной проекта. Допустимая
новизна начинается с единого вывода их сочетания из одной меры и одного
функционала.

\section{Совместная вариация состояния и носителя}

Известны вариационный принцип Гиббса, квантовая относительная энтропия,
метрики Фишера--Боголюбова, спектральное действие и его разложение теплового
ядра, перенормировка и однопетлевые статистические суммы на компактных
седловых геометриях. Проект уже проверял фиксированный
\(\mathbb{RP}^3\times S^1\), энтропийный выбор топологии, нормированные
тепловые состояния, свободную энергию, радиальную вариацию и конечные
контрчлены.

Открытое утверждение уровня \(N_2\): один корреляционный оператор должен
совместно определить состояние и носитель, детерминант флуктуаций, конечные
веса Эйнштейна, квадрата Вейля и Эйлера до сравнения топологий, а затем
пройти физический гессиан и две независимые слепые проверки.

\begin{nogocriterion}[Остановка направления носителя]
Направление закрывается, если порядок топологий зависит от контрчлена,
выбранного после сравнения, или если функционалы Гиббса и положительного
спектрального действия соединяются новым свободным весом.
\end{nogocriterion}

\section{Граничный источник}

Известны спонтанное нарушение \(SO(3)\to A_4\), вильсоновское нарушение
симметрии поколений, динамика Хосотани, конечные вихри порядка три,
канонический фиксированный заряд, роторная эффективная теория и
топологические нулевые моды Майораны. Проект уже построил соответствующие
частные модули, но минимальный механизм Хосотани имеет неверный знак,
вектороподобное исправление не даёт второго сектора, а конечно-геометрическое
действие поколенческого дефекта не проходит знак, идеал лишних форм и
пфаффиан.

Открытое утверждение \(N_2\): одно граничное гильбертово пространство и один
след должны одновременно вывести нечётный фиксированный заряд, ненулевой
конденсат заряда два, физическую тетраэдрическую ось, ровно одну моду
Майораны и вторую чувствительную к Стандартной модели нормировку.

\begin{nogocriterion}[Остановка граничного направления]
Если проектор заряда, кратность поля, ось поколений или второй сектор
задаются независимо после построения, получается собранная вручную модель,
а не общее родительское действие.
\end{nogocriterion}

\section{Конечная геометрия}

Спектральные модели Пати--Салама, обобщённые внутренние флуктуации,
диаграммы Краевского, идеалы лишних форм, вспомогательные поля и
четырёхфермионные каналы от кручения имеют литературные предшественники.
Проект уже проверил ранговые селекторы, внешние квадраты, конусы отображения,
относительные циклы, KO6-дополнения, скрученные связующие поля и высшие
моменты; поколенческая ветвь прошла явное вложение KO6, суперслед,
вспомогательные знаки и гауссов пфаффиан.

Открытое утверждение \(N_2\): ограниченная классификация должна вывести
требуемый относительный знак и невырожденный физический гессиан до
феноменологии.

\begin{nogocriterion}[Остановка конечно-геометрического направления]
Запрещены ещё один проектор ранга один, простое произведение уже проверенных
потенциалов, восстановление лишних форм только на особом фоне и любое
вспомогательное поле без вывода его модуля и меры.
\end{nogocriterion}

\section{Общий вывод и правило дальнейшей работы}

Известны почти все отдельные составляющие. Новизной может быть только их
совместный вывод без свободных коэффициентов. Перед каждой новой проверкой
необходимо указывать ближайшие внутренние предшественники, новые источники
литературы, различать известный элемент, специализацию и новую теорему и
записывать условие остановки. Заявления об абсолютном приоритете запрещены
без независимой проверки специалистом и библиографическими базами.

\begin{protocol}[Литературная разность перед каждой проверкой]
Новая глава допускается лишь после доказательства, что она не переименовывает
закрытый механизм и проверяет одно точно сформулированное совместное
утверждение уровня \(N_2\).
\end{protocol}